\documentclass[../../main.tex]{subfiles}% 注意这里的文件路径不能用 ./main.tex ,否则用latexmk编译子文件会报错
\graphicspath{{\subfix{./image/}}{\subfix{../../image/}}} % 指定图片引用路径,后续可以直接使用图片文件名
% 如果图片存放在上级目录,则使用 ../../image/ 作为备用路径

% 例如:
% \begin{figure}[H]
% \centering
% \includegraphics[width=0.3\textwidth]{图.png}
% \caption{}
% \label{figure:图}
% \end{figure}
% 注意:上述\label{}一定要放在\caption{}之后,否则引用图片序号会只会显示??.

\begin{document}

\section{存在性结果的证明：逼近问题}\label{section:ch07-逼近问题}

设 (Pn) 的初值 $u_{0n}\in L^1(\mathbb{R})\cap L^\infty(\mathbb{R})$ 满足下列性质（见\cref{lemma:ch06-6.5.1}与\eqref{align:ch6-6.49b}--\eqref{align:ch6-6.49c}）：
\begin{align}
u_{0n}\geqslant0\ \text{在 }\mathbb{R},\qquad
\|u_{0n}\|_{L^1(\mathbb{R})}\leqslant\|u_0\|_{\mathscr{R}_f(\mathbb{R})}, \label{align:ch07-7.44}
\end{align}
\begin{align}
u_{0n}\stackrel{*}{\rightharpoonup}u_0,\quad
u_{0n}\to u_{0r}\ \text{在 }\mathbb{R}\text{ 中几乎处处},\quad
\|u_{0n}-u_{0r}\|_{L^1_{\mathrm{loc}}(\mathbb{R}\setminus\operatorname{supp}u_{0s})}\to0. \label{align:ch07-7.45}
\end{align}
设 $\{u_{0n}^\epsilon\}\subseteq C^\infty_c(\mathbb{R})$，$u_{0n}^\epsilon\geqslant0$，为满足
\begin{align}
\|u_{0n}^\epsilon\|_{L^1(\mathbb{R})}\leqslant\|u_{0n}\|_{L^1(\mathbb{R})}\leqslant\|u_0\|_{\mathscr{R}_f(\mathbb{R})},\quad
\|u_{0n}^\epsilon\|_{L^\infty(\mathbb{R})}\leqslant\|u_{0n}\|_{L^\infty(\mathbb{R})}, \label{align:ch07-7.46}
\end{align}
\begin{align}
u_{0n}^\epsilon\to u_{0n}\ \text{在 }L^1(\mathbb{R}),\quad
u_{0n}^\epsilon\stackrel{*}{\rightharpoonup}u_{0n}\ \text{在 }L^\infty(\mathbb{R})\ \text{当 }\epsilon\to0^+. \label{align:ch07-7.47}
\end{align}
设 $\eta\in C^\infty_c(\mathbb{R})$ 为标准磨光核。置 $\eta_\epsilon(y)\triangleq \frac1\epsilon\eta(y/\epsilon)$（$\epsilon>0$），
\begin{align}
\phi_\epsilon(y)\triangleq (\eta_\epsilon*\phi)(y)-(\eta_\epsilon*\phi)(0)\quad(y\in\mathbb{R}), \label{align:ch07-7.48}
\end{align}
其中 $\phi\triangleq \phi\chi_{[0,\infty)}$。与 (Pn) 相伴的正则化问题
\begin{align*}
\begin{cases}
\partial_t u_n^\epsilon+\partial_x[\phi_\epsilon(u_n^\epsilon)]=\epsilon\,\partial_{xx}u_n^\epsilon & \text{在 }S,\\
u_n^\epsilon=u_{0n}^\epsilon & \text{在 }\mathbb{R}\times\{0\},
\end{cases}
\qquad(Pn^\epsilon)
\end{align*}
有唯一强解 $u_n^\epsilon\in C([0,T];H^2(\mathbb{R}))\cap L^\infty(S)$，使 $\partial_t u_n^\epsilon\in L^2(S)$ 且 $u_n^\epsilon\geqslant0$（例如见文献~\cite{ref68}）。

族 $\{u_n^\epsilon\}$ 的若干性质收集于下列引理。第一个引理的证明近乎标准（例如见文献~\cite{ref34}），故从略。

\begin{lemma}\label{lemma:ch07-7.4.1}
设 $u_n^\epsilon$ 为问题 (Pn$^\epsilon$) 的解（$\epsilon>0$，$n\in\mathbb{N}$）。则：
\begin{align}
u_n^\epsilon\geqslant0\ \text{在 }S,\qquad
\|u_n^\epsilon\|_{L^\infty(S)}\leqslant\|u_{0n}\|_{L^\infty(\mathbb{R})}, \label{align:ch07-7.49}
\end{align}
\begin{align}
\|u_n^\epsilon(\cdot,t)\|_{L^1(\mathbb{R})}=\|u_{0n}^\epsilon\|_{L^1(\mathbb{R})}\quad(t\in(0,T)), \label{align:ch07-7.50}
\end{align}
\begin{align}
\|u_n^\epsilon\|_{L^\infty(0,T;L^1(\mathbb{R}))}\leqslant\|u_{0n}\|_{L^1(\mathbb{R})}\leqslant\|u_0\|_{\mathscr{R}_f(\mathbb{R})}, \label{align:ch07-7.51}
\end{align}
\begin{align}
\|u_n^\epsilon(\cdot+h,\cdot)-u_n^\epsilon\|_{L^\infty(0,T;L^1(\mathbb{R}))}
\leqslant\|u_{0n}^\epsilon(\cdot+h)-u_{0n}^\epsilon\|_{L^1(\mathbb{R})}\quad\forall h\in\mathbb{R}. \label{align:ch07-7.52}
\end{align}
\end{lemma}

\begin{lemma}\label{lemma:ch07-7.4.2}
设\eqref{align:ch07-7.1}成立。则存在 $C>0$（仅依赖于 $\|u_0\|_{\mathscr{R}_f(\mathbb{R})}$），使对任意 $n\in\mathbb{N}$、$\epsilon\in(0,1)$ 及 $p\in(0,1)$，
\begin{align}
\epsilon\iint_S(1+u_n^\epsilon)^{p-2}(\partial_x u_n^\epsilon)^2\,\mathrm{d}x\,\mathrm{d}t
\leqslant\frac{C}{p(1-p)}. \label{align:ch07-7.53}
\end{align}
\end{lemma}

\begin{proof}
设 $U\in C^2([0,\infty))$，$U'\geqslant0$ 在 $(0,\infty)$。置
\begin{align}
\Lambda_{U,\epsilon}(y)\triangleq \int_0^y U'(z)\phi_\epsilon'(z)\,\mathrm{d}z+\vartheta_U\quad(\vartheta_U\in\mathbb{R}). \label{align:ch07-7.54}
\end{align}
由\eqref{align:ch07-7.1}与\eqref{align:ch07-7.48}，对所有 $y\geqslant0$，
\begin{align}
|\Lambda_{U,\epsilon}(y)|\leqslant\int_0^y U'(z)|\phi_\epsilon'(z)|\,\mathrm{d}z+|\vartheta_U|
\leqslant M[U(y)-U(0)]+|\vartheta_U|. \label{align:ch07-7.55}
\end{align}
将 (Pn$^\epsilon$) 第一式乘以 $U'(u_n^\epsilon)$，得
\begin{align}
\partial_t[U(u_n^\epsilon)]+\partial_x[\Lambda_{U,\epsilon}(u_n^\epsilon)]
=\epsilon\,\partial_{xx}[U(u_n^\epsilon)]-\epsilon U''(u_n^\epsilon)(\partial_x u_n^\epsilon)^2
\quad\text{在 }S. \label{align:ch07-7.56}
\end{align}
于是对所有 $\zeta\in C^1([0,T];C^2_c(\mathbb{R}))$，
\begin{align}
\begin{aligned}
\epsilon\iint_S&U''(u_n^\epsilon)(\partial_x u_n^\epsilon)^2\zeta\,\mathrm{d}x\,\mathrm{d}t+\int_{\mathbb{R}}U(u_n^\epsilon(x,T))\zeta(x,T)\,\mathrm{d}x\\
&=\int_{\mathbb{R}}U(u_{0n}^\epsilon)\zeta(x,0)\,\mathrm{d}x
+\iint_S[U(u_n^\epsilon)\partial_t\zeta+\Lambda_{U,\epsilon}(u_n^\epsilon)\partial_x\zeta+\epsilon U(u_n^\epsilon)\partial_{xx}\zeta]\,\mathrm{d}x\,\mathrm{d}t.
\end{aligned} \label{align:ch07-7.57}
\end{align}
取 $\vartheta_U=0$，$U(y)=(1+u)^p-1$（$p\in(0,1)$），及
\begin{align*}
\zeta=\rho_k\triangleq \chi_{[-k,k]}+\rho(\cdot-k)\chi_{[k,k+1)}+\rho(\cdot+k)\chi_{(-k-1,-k]}\quad(k\in\mathbb{N})
\end{align*}
其中 $\rho\in C^2_c((-1,1))$ 满足 $\rho(0)=1$、$0\leqslant\rho\leqslant1$、$\rho'(0)=\rho''(0)=0$。则 $0\leqslant U(y)\leqslant y$（$y\geqslant0$），由\eqref{align:ch07-7.51}、\eqref{align:ch07-7.55}与\eqref{align:ch07-7.57}得
\begin{align*}
\epsilon p(1-p)\iint_S(1+u_n^\epsilon)^{p-2}(\partial_x u_n^\epsilon)^2\rho_k\,\mathrm{d}x\,\mathrm{d}t
\leqslant\int_{\mathbb{R}}u_{0n}^\epsilon(x)\,\mathrm{d}x
+\iint_S[Mu_n^\epsilon|\rho_k'|+\epsilon u_n^\epsilon|\rho_k''|]\,\mathrm{d}x\,\mathrm{d}t
\leqslant[1+(M+1)T\|\rho\|_{C^2([-1,1])}]\|u_0\|_{\mathscr{R}_f(\mathbb{R})}=:C
\end{align*}
对所有 $\epsilon\in(0,1)$ 与 $k\in\mathbb{N}$ 成立。令 $k\to\infty$ 得\eqref{align:ch07-7.53}，结论得证。
\end{proof}

\begin{lemma}\label{lemma:ch07-7.4.3}
设\eqref{align:ch07-7.1}成立，$U\in C^2([0,\infty))$。设存在 $K\geqslant0$ 与 $p\in(0,1)$ 使
\begin{align}
|U''(y)|\leqslant K(1+y)^{p-2}\quad\text{对所有 }y\in[0,\infty). \label{align:ch07-7.58}
\end{align}
则存在 $C_p>0$，使对所有 $n\in\mathbb{N}$ 与 $\epsilon>0$，
\begin{align}
\epsilon\iint_S|U''(u_n^\epsilon)|(\partial_x u_n^\epsilon)^2\,\mathrm{d}x\,\mathrm{d}t\leqslant C_p. \label{align:ch07-7.59}
\end{align}
\end{lemma}

\begin{proof}
由\eqref{align:ch07-7.53}与\eqref{align:ch07-7.58}立即得证。
\end{proof}

\begin{lemma}\label{lemma:ch07-7.4.4}
设\eqref{align:ch07-7.1}成立，$U\in C^2([0,\infty))$，$U'\in L^\infty(0,\infty)$ 满足\eqref{align:ch07-7.58}。对任意 $\rho\in C^2_c(\mathbb{R})$ 与 $t\in(0,T)$，置
\begin{align}
U_{n,\rho}^\epsilon(t)\triangleq \int_{\mathbb{R}}U(u_n^\epsilon(x,t))\rho(x)\,\mathrm{d}x,\qquad
U_{n,\rho}(t)\triangleq \int_{\mathbb{R}}U(u_n(x,t))\rho(x)\,\mathrm{d}x. \label{align:ch07-7.60}
\end{align}
则族 $\{U_{n,\rho}^\epsilon\}$ 在 $BV(0,T)$ 中有界。
\end{lemma}

\begin{proof}
由等式\eqref{align:ch07-7.56}，
\begin{align}
(U_{n,\rho}^\epsilon)'(t)=\int_{\mathbb{R}}[\Lambda_{U,\epsilon}(u_n^\epsilon)\rho'+\epsilon U(u_n^\epsilon)\rho''-\epsilon U''(u_n^\epsilon)(\partial_x u_n^\epsilon)^2\rho](x,t)\,\mathrm{d}x. \label{align:ch07-7.61}
\end{align}
因 $U'\in L^\infty(0,\infty)$，存在 $N>0$ 使
\begin{align}
|U(y)|\leqslant N(1+y)\quad\text{对所有 }y\geqslant0. \label{align:ch07-7.62}
\end{align}
故由\eqref{align:ch07-7.1}、\eqref{align:ch07-7.48}与\eqref{align:ch07-7.54}，
\begin{align*}
|\Lambda_{U,\epsilon}(u_n^\epsilon)|\leqslant M\|U'\|_\infty|u_n^\epsilon|+|\vartheta_U|=:\widetilde M u_n^\epsilon+|\vartheta_U|,
\end{align*}
由\eqref{align:ch07-7.61}
\begin{align}
\begin{aligned}
|(U_{n,\rho}^\epsilon)'(t)|&\leqslant\|\rho\|_{C^2(\mathbb{R})}\int_{\operatorname{supp}\rho}\{(\widetilde M+\epsilon N)u_n^\epsilon(x,t)+\epsilon N+|\vartheta_U|\}\,\mathrm{d}x\\
&\quad+\epsilon\|\rho\|_{L^\infty(\mathbb{R})}\int_{\mathbb{R}}|U''(u_n^\epsilon)|(\partial_x u_n^\epsilon)^2(x,t)\,\mathrm{d}x.
\end{aligned} \label{align:ch07-7.63}
\end{align}
由\eqref{align:ch07-7.63}、\eqref{align:ch07-7.51}与\eqref{align:ch07-7.59}，存在 $C_{p,\rho}>0$ 使
\begin{align}
\|(U_{n,\rho}^\epsilon)'\|_{L^1(0,T)}\leqslant\|\rho\|_{C^2(\mathbb{R})}[(\widetilde M+N)T\|u_0\|_{\mathscr{R}_f(\mathbb{R})}+C_{p,\rho}]. \label{align:ch07-7.64}
\end{align}
另一方面，由\eqref{align:ch07-7.51}与\eqref{align:ch07-7.62}，
\begin{align}
\|U_{n,\rho}^\epsilon\|_{L^1(0,T)}\leqslant NT\|\rho\|_{L^\infty(\mathbb{R})}[\|u_0\|_{\mathscr{R}_f(\mathbb{R})}+\lambda(\operatorname{supp}\rho)] \label{align:ch07-7.65}
\end{align}
（$\lambda$ 如惯例表示 $\mathbb{R}$ 中 Lebesgue 测度）。由\eqref{align:ch07-7.64}--\eqref{align:ch07-7.65}得结论。
\end{proof}

置 $R_L\triangleq I_L\times(0,T)$，$I_L\triangleq (-L,L)$。由上述引理得下列收敛结果。

\begin{lemma}\label{lemma:ch07-7.4.5}
设 $\phi\in C([0,\infty))$。则存在子列 $\{u_n^{\epsilon_j}\}\subseteq\{u_n^\epsilon\}$（$\epsilon_j\to0^+$，$j\to\infty$）与 $u_n\in L^\infty(S)\cap L^\infty(0,T;L^1(\mathbb{R}))$、$u_n\geqslant0$，使当 $j\to\infty$ 时
\begin{align}
u_n^{\epsilon_j}\stackrel{*}{\rightharpoonup}u_n\quad\text{在 }L^\infty(S), \label{align:ch07-7.66}
\end{align}
\begin{align}
u_n^{\epsilon_j}(\cdot,t)\stackrel{*}{\rightharpoonup}u_n(\cdot,t)\quad\text{在 }L^\infty(\mathbb{R})\ \text{对几乎所有 }t\in(0,T), \label{align:ch07-7.67}
\end{align}
\begin{align}
u_n^{\epsilon_j}(\cdot,t)\to u_n(\cdot,t)\quad\text{在 }L^1(I_L)\ \text{对几乎所有 }t\in(0,T)\text{ 与所有 }L>0, \label{align:ch07-7.68}
\end{align}
\begin{align}
u_n^{\epsilon_j}\to u_n\quad\text{在 }L^1(R_L)\ \text{对所有 }L>0, \label{align:ch07-7.69}
\end{align}
\begin{align}
u_n^{\epsilon_j}\to u_n,\quad \phi_{\epsilon_j}(u_n^{\epsilon_j})\to\phi(u_n)\quad\text{在 }S\text{ 中几乎处处}. \label{align:ch07-7.70}
\end{align}
此外，
\begin{align}
\|u_n\|_{L^\infty(S)}\leqslant\|u_{0n}\|_{L^\infty(\mathbb{R})},\quad
\|u_n\|_{L^\infty(0,T;L^1(\mathbb{R}))}\leqslant\|u_{0n}\|_{L^1(\mathbb{R})}\leqslant\|u_0\|_{\mathscr{R}_f(\mathbb{R})}. \label{align:ch07-7.71}
\end{align}
\end{lemma}

\begin{proof}
\eqref{align:ch07-7.66}中的收敛、$u_n$ 的非负性及\eqref{align:ch07-7.71}中第一个不等式由\eqref{align:ch07-7.49}立即得到。为证\eqref{align:ch07-7.67}，对任意 $\rho\in C^2_c(\mathbb{R})$ 置
\begin{align*}
I_{n,\rho}^{\epsilon_j}(t)\triangleq \int_{\mathbb{R}}u_n^{\epsilon_j}(x,t)\rho(x)\,\mathrm{d}x\quad(t\in(0,T)).
\end{align*}
由\cref{lemma:ch07-7.4.4}（取 $U(y)=y$），序列 $\{I_{n,\rho}^{\epsilon_j}\}$ 在 $BV(0,T)$ 中有界。故存在子列（为简便不重新标号）与 $I_{n,\rho}\in BV(0,T)$ 使
\begin{align}
I_{n,\rho}^{\epsilon_j}\to I_{n,\rho}\ \text{在 }L^1(0,T)\text{ 且在 }(0,T)\text{ 中几乎处处}. \label{align:ch07-7.72}
\end{align}
另一方面，由\eqref{align:ch07-7.66}，
\begin{align*}
\int_0^T I_{n,\rho}^{\epsilon_j}(t)\,\mathrm{d}t=\iint_S u_n^{\epsilon_j}(x,t)\rho(x)\,\mathrm{d}x\,\mathrm{d}t
\to\iint_S u_n(x,t)\rho(x)\,\mathrm{d}x\,\mathrm{d}t=\int_0^T\mathrm{d}t\int_{\mathbb{R}}u_n(x,t)\rho(x)\,\mathrm{d}x,
\end{align*}
故 $I_{n,\rho}=\int_{\mathbb{R}}u_n(x,t)\rho(x)\,\mathrm{d}x$（对几乎所有 $t\in(0,T)$），且\eqref{align:ch07-7.72}对整个序列成立。综上，对任意 $\rho\in C^2_c(\mathbb{R})$，存在零测集 $N\subseteq(0,T)$ 使
\begin{align*}
\int_{\mathbb{R}}u_n^{\epsilon_j}(x,t)\rho(x)\,\mathrm{d}x\to\int_{\mathbb{R}}u_n(x,t)\rho(x)\,\mathrm{d}x
\end{align*}
对所有 $t\in(0,T)\setminus N$ 成立。因 $C^2_c(\mathbb{R})$ 可分，$N$ 的选取可与 $\rho$ 无关，故\eqref{align:ch07-7.67}成立。

由\eqref{align:ch07-7.51}、\eqref{align:ch07-7.52}与 Fréchet--Kolmogorov 定理，对所有 $t\in(0,T)$ 与 $L>0$，序列 $\{u_n^{\epsilon_j}(\cdot,t)\}$ 在 $L^1(I_L)$ 中相对紧。如前论证，由\eqref{align:ch07-7.67}即得\eqref{align:ch07-7.68}。由\eqref{align:ch07-7.51}与\eqref{align:ch07-7.68}得
\begin{align*}
\|u_n\|_{L^\infty(0,T;L^1(I_L))}\leqslant\|u_{0n}\|_{L^1(\mathbb{R})}\leqslant\|u_0\|_{\mathscr{R}_f(\mathbb{R})}\quad\forall L>0,
\end{align*}
由 $L$ 任意得\eqref{align:ch07-7.71}中第二个不等式。由\eqref{align:ch07-7.51}与\eqref{align:ch07-7.68}及控制收敛定理得\eqref{align:ch07-7.69}中收敛。最后，\eqref{align:ch07-7.70}中第一个收敛由\eqref{align:ch07-7.69}（可能抽取子列，不重新标号）得到。因对任意紧集 $K\subseteq\mathbb{R}$，当 $\epsilon\to0^+$ 时 $\|\phi_\epsilon-\phi\|_{C(K)}\to0$，\eqref{align:ch07-7.70}中第二个收敛亦得证。
\end{proof}

\begin{lemma}\label{lemma:ch07-7.4.6}
设\eqref{align:ch07-7.1}成立，$U\in C^2([0,\infty))$，$U'\in L^\infty(0,\infty)$ 满足\eqref{align:ch07-7.58}。对任意 $\rho\in C^2_c(\mathbb{R})$ 与 $t\in(0,T)$，置
\begin{align}
U_{n,\rho}(t)\triangleq \int_{\mathbb{R}}U(u_n(x,t))\rho(x)\,\mathrm{d}x \label{align:ch07-7.73}
\end{align}
其中 $u_n\in L^\infty(S)\cap L^\infty(0,T;L^1(\mathbb{R}))$ 由\cref{lemma:ch07-7.4.5}给出。则：
\begin{enumerate}[(i)]
\item 存在子列 $\{U_{n,\rho}^{\epsilon_j}\}\subseteq\{U_{n,\rho}^\epsilon\}$（$\epsilon_j\to0^+$）（$U_{n,\rho}^\epsilon$ 由\eqref{align:ch07-7.60}定义）使
\begin{align}
U_{n,\rho}^{\epsilon_j}\to U_{n,\rho}\quad\text{在 }L^1(0,T); \label{align:ch07-7.74}
\end{align}
\item 族 $\{U_{n,\rho}\}$ 在 $BV(0,T)$ 中有界。
\end{enumerate}
\end{lemma}

\begin{proof}
关于 (i)，设 $\{u_n^{\epsilon_j}\}$ 为\cref{lemma:ch07-7.4.5}中子列。则对几乎所有 $t\in(0,T)$，由\eqref{align:ch07-7.70}中第一个收敛有 $U(u_n^{\epsilon_j}(\cdot,t))\to U(u_n(\cdot,t))$ 在 $\mathbb{R}$ 中几乎处处，由控制收敛定理得\eqref{align:ch07-7.74}。关于 (ii)，由\eqref{align:ch07-7.74}与\eqref{align:ch07-7.65}得
\begin{align}
\|U_{n,\rho}\|_{L^1(0,T)}=\lim_{j\to\infty}\|U_{n,\rho}^{\epsilon_j}\|_{L^1(0,T)}
\leqslant NT\|\rho\|_{L^\infty(\mathbb{R})}[\|u_0\|_{\mathscr{R}_f(\mathbb{R})}+\lambda(\operatorname{supp}\rho)]. \label{align:ch07-7.75}
\end{align}
另一方面，由\eqref{align:ch07-7.64}与 $L^1(0,T)$ 拓扑下总变差的下半连续性（见\cref{remark:ch03-3.1.4}(ii)），
\begin{align}
\|U_{n,\rho}'\|_{\mathscr{R}_f(0,T)}\leqslant\|\rho\|_{C^2(\mathbb{R})}[(\widetilde M+N)T\|u_0\|_{\mathscr{R}_f(\mathbb{R})}+C_{p,\rho}]. \label{align:ch07-7.76}
\end{align}
由\eqref{align:ch07-7.75}--\eqref{align:ch07-7.76}得 (ii)。
\end{proof}

回忆下列定义（例如见文献~\cite{ref34}）。

\begin{definition}\label{definition:ch07-Pn熵解}
函数 $u_n\in L^\infty(0,T;L^1(\mathbb{R}))\cap L^\infty(S)$ 称为问题 (Pn) 的熵解，若对每个 $\zeta\in C^1([0,T];C^1_c(\mathbb{R}))$、$\zeta(\cdot,T)=0$ 在 $\mathbb{R}$ 上、$\zeta\geqslant0$ 及任意 $l\in[0,\infty)$，有
\begin{align}
\iint_S[E_l(u_n)\partial_t\zeta+F_l(u_n)\partial_x\zeta]\,\mathrm{d}x\,\mathrm{d}t
+\int_{\mathbb{R}}E_l(u_{0n})\zeta(x,0)\,\mathrm{d}x\geqslant0, \label{align:ch07-7.77}
\end{align}
其中 $E_l(y)=|y-l|$，$F_l(y)=\operatorname{sgn}(y-l)[\phi(y)-\phi(l)]$（$y\in[0,\infty)$）。
\end{definition}

\begin{remark}\label{remark:ch07-7.4.1}
熵解是弱解：若 $\zeta\in C^1([0,T];C^1_c(\mathbb{R}))$、$\zeta(\cdot,T)=0$ 在 $\mathbb{R}$ 上，则
\begin{align}
\iint_S[u_n\partial_t\zeta+\phi(u_n)\partial_x\zeta]\,\mathrm{d}x\,\mathrm{d}t
+\int_{\mathbb{R}}u_{0n}\zeta(x,0)\,\mathrm{d}x=0. \label{align:ch07-7.78}
\end{align}
\end{remark}

\begin{proposition}\label{proposition:ch07-7.4.7}
设 $\phi\in C([0,\infty))$。则\cref{lemma:ch07-7.4.5}给出的函数 $u_n$ 是问题 (Pn)（$n\in\mathbb{N}$）的熵解。此外，
\begin{align}
\|u_n(\cdot,t)\|_{L^1(\mathbb{R})}=\|u_{0n}\|_{L^1(\mathbb{R})}\quad\text{对几乎所有 }t\in(0,T) \label{align:ch07-7.79}
\end{align}
且
\begin{align}
\|u_n(\cdot+h,\cdot)-u_n\|_{L^\infty(0,T;L^1(\mathbb{R}))}
\leqslant\|u_{0n}(\cdot+h)-u_{0n}\|_{L^1(\mathbb{R})}\quad\forall h\in\mathbb{R}. \label{align:ch07-7.80}
\end{align}
若 $\phi$ 局部 Lipschitz 连续，则 $u_n$ 是 (Pn) 的唯一熵解。
\end{proposition}

\begin{proof}
设 $\zeta$ 如\cref{definition:ch07-Pn熵解}，$E_l(y)=|y-l|$，$F_{l,\epsilon}(y)=\operatorname{sgn}(y-l)[\phi_\epsilon(y)-\phi_\epsilon(l)]$（$y,l\in[0,\infty)$）。由 (Pn$^\epsilon$) 作标准计算得
\begin{align}
\iint_S\{E_l(u_n^\epsilon)(\partial_t\zeta+\epsilon\zeta_{xx})+F_{l,\epsilon}(u_n^\epsilon)\partial_x\zeta\}\,\mathrm{d}x\,\mathrm{d}t
+\int_{\mathbb{R}}E_l(u_{0n}^\epsilon)\zeta(x,0)\,\mathrm{d}x\geqslant0. \label{align:ch07-7.81}
\end{align}
设 $\{u_n^{\epsilon_j}\}\subseteq\{u_n^\epsilon\}$ 由\cref{lemma:ch07-7.4.5}给出。由\eqref{align:ch07-7.48}--\eqref{align:ch07-7.49}，
\begin{align*}
\|\phi_{\epsilon_j}(u_n^{\epsilon_j})\|_{L^\infty(S)}
\leqslant\sup_{|y|\leqslant\|u_{0n}\|_{L^\infty(\mathbb{R})}}|\phi_{\epsilon_j}(y)|
\leqslant\sup_{|y|\leqslant\|u_{0n}\|_{L^\infty(\mathbb{R})}}|\phi(y)|.
\end{align*}
由\eqref{align:ch07-7.70}、上式与控制收敛定理，
\begin{align*}
\iint_S F_{l,\epsilon_j}(u_n^{\epsilon_j})\partial_x\zeta\,\mathrm{d}x\,\mathrm{d}t
\to\iint_S F_l(u_n)\partial_x\zeta\,\mathrm{d}x\,\mathrm{d}t.
\end{align*}
由\eqref{align:ch07-7.69}与\eqref{align:ch07-7.47}亦有
\begin{align*}
\iint_S E_l(u_n^{\epsilon_j})(\partial_t\zeta+\epsilon_j\zeta_{xx})\,\mathrm{d}x\,\mathrm{d}t
\to\iint_S E_l(u_n)\partial_t\zeta\,\mathrm{d}x\,\mathrm{d}t,
\end{align*}
\begin{align*}
\int_{\mathbb{R}}E_l(u_{0n}^{\epsilon_j})\zeta(x,0)\,\mathrm{d}x\to\int_{\mathbb{R}}E_l(u_{0n})\zeta(x,0)\,\mathrm{d}x.
\end{align*}
在\eqref{align:ch07-7.81}（取 $\epsilon=\epsilon_j$）中令 $j\to\infty$ 得\eqref{align:ch07-7.77}，故 $u_n$ 是 (Pn) 的熵解。唯一性由 Kružkov 唯一性定理得到（例如见文献~\cite{ref91}）。

不等式\eqref{align:ch07-7.80}由\eqref{align:ch07-7.52}与\eqref{align:ch07-7.68}得到。关于\eqref{align:ch07-7.79}，由\eqref{align:ch07-7.78}经标准论证得，对几乎所有 $t\in(0,T)$ 与所有 $\rho\in C^1_c(\mathbb{R})$，
\begin{align}
\int_{\mathbb{R}}u_n(x,t)\rho(x)\,\mathrm{d}x-\int_{\mathbb{R}}u_{0n}(x)\rho(x)\,\mathrm{d}x
=\int_0^t\int_{\mathbb{R}}\phi(u_n)(x,s)\rho'(x)\,\mathrm{d}x\,\mathrm{d}s. \label{align:ch07-7.82}
\end{align}
设 $\rho_k\in C^1_c(\mathbb{R})$（$k\in\mathbb{N}$）满足
\begin{align*}
\rho_k(x)=\begin{cases}1&\text{若 }|x|\leqslant k,\\0&\text{若 }|x|\geqslant k+1,\end{cases}
\end{align*}
且 $\|\rho_k'\|_{L^\infty(\mathbb{R})}\leqslant2$。因 $u_n\in L^1(S)$，在\eqref{align:ch07-7.82}中取 $\rho=\rho_k$ 并令 $k\to\infty$，利用\eqref{align:ch07-7.1}得
\begin{align*}
\left|\int_0^t\int_{\mathbb{R}}\phi(u_n)(x,s)\rho_k'(x)\,\mathrm{d}x\,\mathrm{d}s\right|
\leqslant2M\int_0^t\int_{\{k\leqslant|x|\leqslant k+1\}}|u_n(x,s)|\,\mathrm{d}x\,\mathrm{d}s\to0.
\end{align*}
另一方面，由单调收敛定理，
\begin{align*}
\int_{\mathbb{R}}u_n(x,t)\rho_k(x)\,\mathrm{d}x\to\int_{\mathbb{R}}u_n(x,t)\,\mathrm{d}x,\quad
\int_{\mathbb{R}}u_{0n}(x)\rho_k(x)\,\mathrm{d}x\to\int_{\mathbb{R}}u_{0n}(x)\,\mathrm{d}x.
\end{align*}
由\eqref{align:ch07-7.82}得等式\eqref{align:ch07-7.79}。
\end{proof}

为后文引用，证明下列结果。

\begin{lemma}\label{lemma:ch07-7.4.8}
设\eqref{align:ch07-7.1}成立，$U\in C^2([0,\infty))\cap L^\infty(0,\infty)$ 满足\eqref{align:ch07-7.58}。置
\begin{align}
\Lambda_U(y)\triangleq \int_0^y U'(z)\phi'(z)\,\mathrm{d}z+\vartheta_U\quad(y,\vartheta_U\in\mathbb{R}), \label{align:ch07-7.83}
\end{align}
并设 $\Lambda_U\in L^\infty(0,\infty)$。设 $u_n$ 为由\cref{proposition:ch07-7.4.7}给出的问题 (Pn)（$n\in\mathbb{N}$）的熵解。则对任意 $n\in\mathbb{N}$，存在 $\sigma_n\in\mathscr{R}_f(S)$，使对所有 $\zeta\in C^1_c(S)$，
\begin{align}
\iint_S\{U(u_n)\partial_t\zeta+\Lambda_U(u_n)\partial_x\zeta\}\,\mathrm{d}x\,\mathrm{d}t=\langle\sigma_n,\zeta\rangle. \label{align:ch07-7.84}
\end{align}
此外有
\begin{align}
\|\sigma_n\|_{\mathscr{R}_f(S)}\leqslant C_p \label{align:ch07-7.85}
\end{align}
其中 $C_p>0$ 如\eqref{align:ch07-7.59}。
\end{lemma}

\begin{proof}
由不等式\eqref{align:ch07-7.59}，对任意固定的 $n\in\mathbb{N}$，族 $\{\epsilon U''(u_n^\epsilon)(\partial_x u_n^\epsilon)^2\}$ 在 $L^1(S)\subseteq\mathscr{R}_f(S)$ 中有界。故由 Banach--Alaoglu 定理，存在序列 $\{\epsilon_k\}$（$\epsilon_k\to0^+$，$k\to\infty$）与 $\sigma_n\in\mathscr{R}_f(S)$ 使
\begin{align}
\epsilon_k U''(u_n^{\epsilon_k})(\partial_x u_n^{\epsilon_k})^2\stackrel{*}{\rightharpoonup}\sigma_n\quad\text{当 }k\to\infty. \label{align:ch07-7.86}
\end{align}
由范数的下半连续性（见\cref{remark:ch05-5.1.3}），从\eqref{align:ch07-7.86}得不等式\eqref{align:ch07-7.85}。

设 $\zeta\in C^2_c(S)$。由\eqref{align:ch07-7.56}，
\begin{align}
\epsilon\iint_S U''(u_n^\epsilon)(\partial_x u_n^\epsilon)^2\zeta\,\mathrm{d}x\,\mathrm{d}t
=\iint_S[U(u_n^\epsilon)\partial_t\zeta+\Lambda_{U,\epsilon}(u_n^\epsilon)\partial_x\zeta+\epsilon U(u_n^\epsilon)\partial_{xx}\zeta]\,\mathrm{d}x\,\mathrm{d}t \label{align:ch07-7.87}
\end{align}
其中 $\Lambda_{U,\epsilon}$ 由\eqref{align:ch07-7.54}定义。由\eqref{align:ch07-7.1}与\eqref{align:ch07-7.49}中第二个不等式，对任意固定的 $n\in\mathbb{N}$，
\begin{align*}
|\Lambda_{U,\epsilon}(u_n^\epsilon)|\leqslant M\int_0^{\|u_{0n}\|_{L^\infty(\mathbb{R})}}|U'(z)|\,\mathrm{d}z+|\vartheta_U|<\infty,
\end{align*}
故对固定的 $n$，族 $\{\Lambda_{U,\epsilon}(u_n^\epsilon)\}$ 在 $L^\infty(S)$ 中有界。在\eqref{align:ch07-7.87}（取 $\epsilon=\epsilon_k$）中令 $k\to\infty$，如\cref{proposition:ch07-7.4.7}证明中那样论证，即得等式\eqref{align:ch07-7.84}。
\end{proof}




\end{document}